\section{A new source invariant and two regions of unreachable code}
\label{sec:source-product-invariant}

The following result concerns the exact 6,455-leaf source program of the
237-state checkpoint. It is an ordinary inductive proof supported
by a finite certificate, and it does not itself reduce the binary state count.
It supplies two precise regions in which source code may be changed without
affecting the initialized proof-search computation.

The source list has SHA-256
\begin{quote}\small\ttfamily
7fd152d1f5371d2a561d4ad1870d241683\allowbreak
b010ce7b7a27399b95363a4537cfb6.
\end{quote}
Its ordered instruction indices start at zero.  Registers are logical
registers, before the inherited physical permutation.  The initial register
file is
\[
(3,0,0,0,3,0,0,0,0).
\]
The finite graph interface is the inherited ordered-leaf graph: an increment
advances the program counter, a positive conditional decrement decreases its
register and advances, and a zero conditional decrement executes a level-one
break.  The graph extraction checks the required zero-padded path convention.

\subsection{A finite abstraction with an unbounded meaning}

Use the eighteen classes
\[
\{0\},\{1\},\ldots,\{15\},\qquad
 E=\{16,18,20,\ldots\},\qquad
 O=\{17,19,21,\ldots\}.
\]
An abstract register value is any union of these classes.  An abstract register
file is the Cartesian product of nine such unions.  Different program points
have different abstract register files; an unreached point is assigned the
empty set.  Class unions are encoded as eighteen-bit nonnegative integers.

Increment is exact on the singleton classes below 15, maps 15 into $E$, maps
$E$ into $O$, and maps $O$ into $E$.  Positive decrement is exact on positive
singletons, and its two large-class images are
\[
\operatorname{pred}(E)=\{15\}\cup O,
\qquad
\operatorname{pred}(O)=E.
\]
Transfer across a zero decrement branch restricts its tested register to
$\{0\}$; transfer across a positive branch first excludes zero and then applies
positive decrement.  Jumps preserve all register values.  Every abstract image described here contains the concrete image.  The
displayed predecessor images of $E$ and $O$ are exact; mapping 15 into $E$
and incrementing $O$ can overapproximate the concrete image.  Taking Cartesian
products can additionally discard relations between registers, but it only
enlarges the possible set of files.

Starting with the exact initialized seed at the first decoded instruction,
the replay iterates these transfers and componentwise union until a fixed
point is reached.  The lattice is finite, so this process terminates without
a bound on the source computation, register values, or the number of proof
search iterations.  For the threshold-16 run, 273,541 strict point updates
produce the retained certificate.  A separate pass checks every applicable
edge against the final annotation.

\paragraph{Inductive invariant theorem.}
Every initialized source configuration belongs to the abstract register file
assigned to its decoded instruction point.

\paragraph{Proof.}
The initial seed is included in its point annotation.  Each concrete source
step is included in the corresponding abstract transfer just described.  The
post-fixpoint condition says that the resulting abstract file is included in
the annotation at its actual successor.  Induction on the number of source
steps proves the assertion.  The induction ranges over arbitrary natural
step counts.  All underlying objects are finite strings, natural numbers, and
finite lists; the argument is formalizable in ZFC.  The Python replay is a
mathematical checker, not a new kernel-checked theorem.

\subsection{Exact consequences}

The resulting annotation reaches 6,452 instruction points and excludes exactly
three.  It proves the following stronger zero-register facts:
\[
\begin{array}{c|l|l}
\text{index}&\text{instruction path}&\text{registers forced to zero}\\\hline
6314&\texttt{111010010110001110}&R_3,R_6,R_7\\
6429&\texttt{111111000}&R_1,R_3,R_5,R_6,R_7
\end{array}
\]
The instructions at these two points are respectively $\operatorname{CD}(6)$
and $\operatorname{CD}(1)$.  Their positive branches are therefore impossible
in the initialized source computation.  The three excluded points are
\[
\begin{array}{c|l|l}
\text{index}&\text{path}&\text{instruction}\\\hline
6315&\texttt{111010010110001111}&\operatorname{CC}(2)\\
6430&\texttt{1111110010}&\operatorname{CI}(6)\\
6431&\texttt{1111110011}&\operatorname{CC}(2).
\end{array}
\]
Thus the second and third excluded leaves form one wholly unreachable subtree
at prefix \texttt{111111001}.  A simpler three-class abstraction using zero,
positive even, and odd already proves the last two leaves unreachable and the
zero fact at point 6429.  Increasing the exact-value threshold from 16 to 64
produces the same three excluded points and two forced decrement guards; this
is a scoped observation, not a completeness result for source invariants.

\subsection{A replacement theorem for the initialized computation}

At path \texttt{111010010110001110}, replace $\operatorname{CD}(6)$ by any of
\[
\operatorname{CB}(1),\quad
\operatorname{CD}(3),\quad \operatorname{CD}(6),\quad \operatorname{CD}(7).
\]
At path \texttt{111111000}, replace $\operatorname{CD}(1)$ by any of
\[
\operatorname{CB}(1),\quad
\operatorname{CD}(1),\quad \operatorname{CD}(3),\quad
\operatorname{CD}(5),\quad \operatorname{CD}(6),\quad \operatorname{CD}(7).
\]
Finally replace the subtree at either or both prefixes
\[
\texttt{111010010110001111},\qquad \texttt{111111001}
\]
by any finite source subtrees, leaving their attachment prefixes and all other
source structure fixed.  No remaining reachable instruction path is changed.

\paragraph{Source replacement theorem.}
Every such replacement gives exactly the same sequence of register files and
source program-counter words as the original program from the displayed seed.
In particular its source halting behavior is identical.

\paragraph{Proof.}
Compare the two source computations by induction.  The original computation
never decodes a point inside either replaced subtree, by the invariant.  At a
retained forced-zero guard, every permitted replacement is a zero conditional
decrement or an explicit level-one break.  Each preserves the entire register
file and applies the identical level-one-break operation to the identical
program-counter word.  Every other decoded instruction is unchanged.
Consequently one old step and one new step have exactly the same successor
configuration.  The induction keeps the two trajectories equal and hence also
keeps the new trajectory away from the replaced subtrees.  This proves both
halting and nonhalting equivalence without any consistency assumption.

This theorem is specific to the initialized seed.  It makes no claim about
arbitrary register files at the affected instructions.  A compiler or binary
runtime proof that requires every syntactic instruction guard must either
retain that broader interface or explicitly incorporate this new invariant.
The inherited ZFC calibration can be composed with the source equivalence
once a complete binary simulation and a correct blank initializer for a new
table have been established.

\subsection{A checked contraction and its present limitation}

As an additional experiment, remove the entire zero loop at prefix
\texttt{11111100}, consisting of old points 6429--6431, and contract its
surviving sibling \texttt{11111101} into \texttt{1111110}.  The resulting source
has 6,452 leaves.  Every surviving instruction label and jump operand remains
unchanged, and all its successor indices agree with the old graph after
sending the removed zero loop to old point 6432.  The original computation
can make only one consecutive step inside the removed loop, so deleting that
step is a proper stuttering simulation and preserves infinite computations.

However, hash-consing of the source decision tree, treating each complete
instruction as an atomic leaf, changes its number of unique internal subtrees
from 201 to 202.  Therefore this contraction does not
supply a smaller binary machine.  The more useful optimization interface is
the preceding replacement theorem, which preserves all reachable path lengths
and permits the unreachable regions to be chosen for decoder sharing.  No
236-state table or other sub-237-state binary result is claimed by this
invariant package.  Initializer correctness must be checked independently for
any later literal decoder change, because the initializer shares states with
the running program.

\subsection{Independent verification and relational refinement}

A second checker imports no discovery code. It reconstructs the binary trie and the literal consumed-path control operations, then verifies seed inclusion and all 8,380 enabled nonhalting edges. It rejects both excluding a seed value and admitting a positive value at a forced-zero guard. The retained receipt is \file{evidence/source_product_k16_independent.json}.

The relational extension tracks an arbitrary set of the 512 joint parity patterns of nine registers and reduces it against the product value classes after every transfer and join. Its final post-fixpoint scan passes, but it proves the same two forced-zero guards and three unreachable leaves. This is a finite abstraction of an unbounded execution, and its negative discovery result is limited to the chosen domain.

\subsection{Instruction contracts restricted by the source invariant}

The source invariant can justify a stronger binary optimization interface than
testing every instruction on every register file.  Let $I_i\subseteq\mathbb N^9$
be the product of the certified register sets at leaf $i$, with $I_i=\varnothing$
for a bottom fact.  The independent checker reconstructs the complete source
control graph from a full binary trie and the consumed-path list operations.
It verifies the initial seed and every guarded transfer inclusion.  Therefore
every finite prefix of the actual source computation reaches a register tuple
in the corresponding $I_i$.

For the delivered threshold-$16$ certificate, the abstract atoms are
\[
 \{0\},\{1\},\ldots,\{15\},\qquad
 E=\{16,18,20,\ldots\},\qquad
 O=\{17,19,21,\ldots\}.
\]
The threshold is not a bound on any register.  The sets $E$ and $O$ retain
arbitrarily large values.  Their arithmetic images follow from elementary
progression identities:
\[
 E+1=O,\quad O+1=E\setminus\{16\},\quad
 E-1=\{15\}\cup O,\quad O-1=E.
\]
The least abstract image of $O+1$ is $E$, so the transfer deliberately
overapproximates that concrete image.  Singleton images are handled exactly
until they enter a tail atom.  Taking unions proves all $36$ atomic increment
and decrement rules.  No numerical sample is substituted for these identities.

The independent check closes $8{,}380$ enabled nonhalting edges and the single
halt leaf.  It proves that logical register $R_6$ is zero at leaf $6314$, whose
path is \texttt{111010010110001110}, and that $R_1$ is zero at leaf $6429$,
whose path is \texttt{111111000}.  Leaves $6315$, $6430$, and $6431$ are
unreachable.  These facts concern the calibrated source computation from its
specified seed.  They do not assert the same restrictions on arbitrary
register frames supplied externally to the binary interpreter.

\paragraph{Affine covering families.}
The literal unary/list engine accepts affine natural-valued register
expressions.  To use it without splitting a product into exponentially many
atomic cases, write one coordinate of $I_i$ as a union of a finite set $F$ of
small values and zero, one, or two progressions with first terms in $B$ and
common step $2$.  Let
\[
 m=\min(F\cup B),\qquad
 d=\gcd\Bigl(\{a-m:a\in F\cup B\}
       \cup\{2:B\ne\varnothing\}\Bigr),
\]
where the gcd of a set consisting only of zero is zero.  The notation for the
second set means that $2$ is included precisely when a tail is present.
Every allowed value belongs to $m+d\mathbb N$; if $d=0$, this family is the
singleton $\{m\}$.  Indeed, each finite difference is divisible by $d$, and
for a tail value $a+2t$ both $a-m$ and $2$ are divisible by $d$.

Before forming these families for a decrement, intersect the selected
coordinate with its zero or positive guard.  A zero branch thus receives the
exact constant zero; every positive branch has $m\ge1$.  Empty branches and
bottom leaves are omitted only because the independent source invariant
proves them impossible.  Each other coordinate receives its own independent
natural variable.  The gap, first spare, and arbitrary finite spare tail remain
unrestricted.  This leaves $8{,}381$ universal source interfaces, including the
halt interface, compared with the original $8{,}386$ unrestricted interfaces.

The wrapper \file{verification/check_source_restricted_runtime.py} verifies
the source invariant first, checks that the runtime helper's path arithmetic
matches the independently reconstructed control graph, and then substitutes
these affine families into the unchanged three-rule execution engine.  It
rederives the four finite-list schemas from the supplied literal table.  Each
nonhalting interface must still reach its complete prescribed tape, state,
orientation, and head displacement with positive finite duration; a halt
interface must reach the actual halt state.  The output retains every family,
guard, affine clock, exact input digest, and complete case coverage.

\begin{proposition}[Source-restricted binary simulation]
Suppose the delivered source invariant contains the calibrated seed and is
preserved by every source step.  Suppose a binary table has positive finite
instruction contracts for affine families covering every permitted guarded
product $I_i$, with exact source updates and whole-tape endpoints.  If its
blank initializer reaches the prescribed source frame, then its blank halting
is equivalent to halting of that calibrated source computation.
\end{proposition}

\begin{proof}
At the initial frame the actual source registers belong to $I_0$.  Assume the
source and binary computations agree at a frame for leaf $i$.  The actual
tuple lies in $I_i$, so it lies in one of the permitted guarded affine
families.  The corresponding positive binary block reaches exactly the frame
of the source successor.  Source-invariant preservation places the new actual
tuple in the successor product.  Induction gives corresponding frames for
every source prefix.

A source halt supplies a finite binary halting block.  If the source run is
infinite, the positive nonhalting blocks concatenate into an infinite binary
run.  Such a block cannot hide an interior halt because the machine has no
transition out of the halt outcome.  The two computations therefore halt
together.  Notice that the enlarged affine families themselves need not be
closed under source transitions: the induction uses the narrower certified
products for the actual run.  All arguments are ordinary finite syntax,
natural-number arithmetic, and induction, and are formalizable in ZFC.
\end{proof}

The initializer premise remains explicit.  A successful restricted instruction
replay alone does not certify a new machine's blank startup, and unreachable
source leaves alone do not remove literal binary states from a shared decoder.

The complete fresh restricted replay succeeds on all 8,381 interfaces with 2,249,105 symbolic rules and minimum duration 13. It still directly uses all 474 rows of the 237-state table, so these reachability facts do not expose an entirely unused binary state. The complete case-array SHA-256 is
\hash{b359d39135f4cb12323b28a05ea56513}{14a96d10afc8da52dfa3bd5ff8741e9d}
